%%%%%%%%%%%%%%%%%%%%%%%%%%%%%%%%%%%%%%%
% Cover Letter - AirHelp, Lead Product Designer (Mobile App), Berlin
%%%%%%%%%%%%%%%%%%%%%%%%%%%%%%%%%%%%%%%

\documentclass[]{cover-fira}
\usepackage{fancyhdr}

\pagestyle{fancy}
\fancyhf{}

\rfoot{Page \thepage \hspace{0pt}}
\thispagestyle{empty}
\renewcommand{\headrulewidth}{0pt}
\begin{document}

%%%%%%%%%%%%%%%%%%%%%%%%%%%%%%%%%%%%%%
%     TITLE NAME
%%%%%%%%%%%%%%%%%%%%%%%%%%%%%%%%%%%%%%
\namesection{}{\Huge{Marko Rosić}}{%
  \href{mailto:marko@rosic.net}{marko@rosic.net} | +381 60 585 44 55 |
  \urlstyle{same}\href{https://linkedin.com/in/markorosic}{LinkedIn}
}

%%%%%%%%%%%%%%%%%%%%%%%%%%%%%%%%%%%%%%
%     MAIN COVER LETTER CONTENT
%%%%%%%%%%%%%%%%%%%%%%%%%%%%%%%%%%%%%%

\currentdate{\today}
\lettercontent{Dear AirHelp hiring team,}

\lettercontent{I am applying for the Lead Product Designer role on the AirHelp mobile app. Flight compensation is a high-stakes, time-sensitive journey: passengers are frustrated, the process is opaque, and the moment they open the app is rarely a calm one. Designing that experience well — clearly, quickly, with the right emotional register — is exactly the kind of problem I find most interesting.}

\lettercontent{Two things about this role stood out: it is hypothesis-driven ("form hypotheses, build and test solutions") and it explicitly names Claude as a prototyping tool. That is how I work. I use Claude Code daily to prototype flows and build tooling, which lets me validate ideas before committing to a full design direction. My background:}

{\raggedright\fontspec{Fira Sans}\fontsize{11pt}{13pt}\selectfont
\begin{itemize}
    \item \textbf{Hypothesis-driven design at scale:} at Catena Media I ran A/B testing, analytics review, and user research as a continuous loop — not a phase. Every redesign decision was framed as a hypothesis and measured against traffic and retention outcomes. Mobile traffic rose 20\%, retention 15\%.
    \item \textbf{Design systems that hold up:} I build token-based component libraries that engineering teams actually use. One cut development time by 30\% and scaled across 10+ websites; another cut time-to-market by 25\% across a multi-brand portfolio.
    \item \textbf{Mobile product experience:} at Smith Micro I designed carrier-branded Android apps shipped on Sprint and T-Mobile devices, working inside strict brand and device constraints across multiple simultaneous variants.
    \item \textbf{AI-accelerated prototyping:} the JD names Claude — I use it to go from idea to working prototype in an afternoon. It changes the speed at which hypotheses can be tested, which is the part of the design process I care most about.
\end{itemize}\par}
\vspace{6pt}

\lettercontent{AirHelp operates across multiple markets with millions of passengers annually. The app is the primary touchpoint at a moment of genuine user need. I would like to bring the same rigour and speed I have applied to high-traffic consumer products to that challenge.}

\lettercontent{I am based in Belgrade (CET) and open to relocation to Berlin. I would be glad to talk about the role.}

\begin{flushright}
\closing{Kind regards,}

\signature{Marko Rosić}
\end{flushright}
\end{document}
